\documentclass{article}
\usepackage{amsmath,amssymb,amsthm}
\usepackage{algorithm}
\usepackage{algpseudocode}
\usepackage{geometry}
\geometry{margin=1in}
\newtheorem{theorem}{Theorem}
\newtheorem{lemma}{Lemma}
\newtheorem{assumption}{Assumption}
\newtheorem{remark}{Remark}
\newcommand{\grad}{\nabla}
\newcommand{\E}{\mathbb{E}}

\begin{document}

%%%%%%%%%%%%%%%%%%%%%%%%%%%%%%%%%%%%%%%%%%%%%%%%%%%%%%%%%%%%
\section*{Heavy–Ball Method (Momentum Gradient Descent)}
%%%%%%%%%%%%%%%%%%%%%%%%%%%%%%%%%%%%%%%%%%%%%%%%%%%%%%%%%%%%

\section*{Mathematical Formulation}
Let $f:\mathbb{R}^d\to\mathbb{R}$ be differentiable.  
Given an initial pair $(x_0,x_{-1})\in\mathbb{R}^d\times\mathbb{R}^d$, a stepsize $\alpha>0$ and a momentum parameter $\beta\in[0,1)$, the Heavy–Ball iterates are defined for $k\ge 0$ by
\begin{equation}
\boxed{
x_{k+1}=x_k+\beta\,(x_k-x_{k-1})-\alpha\,\grad f(x_k).
}
\label{eq:HB}
\end{equation}
The term $\beta(x_k-x_{k-1})$ is a \emph{velocity} built from the previous displacement; the gradient is evaluated at the \emph{current} point $x_k$ (no look‑ahead).

\section*{Pseudocode}
\begin{algorithm}[H]
\caption{Heavy–Ball Method}
\begin{algorithmic}[1]
\State \textbf{Input:} stepsize $\alpha>0$, momentum $\beta\in[0,1)$, tolerance $\varepsilon>0$, max iter $K$.
\State \textbf{Initialize:} $x_0\in\mathbb{R}^d$, $x_{-1}=x_0$.
\For{$k=0,1,\dots,K-1$}
    \State $g_k\gets \grad f(x_k)$
    \State $x_{k+1}\gets x_k+\beta (x_k-x_{k-1})-\alpha g_k$
    \If{$\|g_k\|\le\varepsilon$}
        \State \textbf{break}
    \EndIf
\EndFor
\State \textbf{Output:} $x_{k+1}$
\end{algorithmic}
\end{algorithm}

\section*{Dry Run Example}
Consider the one–dimensional quadratic
\[
f(w)=(w-3)^2,
\qquad \grad f(w)=2(w-3).
\]
Set $\alpha=0.1$, $\beta=0.5$, $w_{-1}=w_0=0$.

\begin{center}
\begin{tabular}{c|c|c|c}
Iteration $k$ & $g_k=2(w_k-3)$ & $w_{k+1}=w_k+\beta(w_k-w_{k-1})-\alpha g_k$ & $f(w_{k+1})$ \\ \hline
0 & $2(0-3)=-6$ & $0+0.5(0-0)-0.1(-6)=0.6$ & $(0.6-3)^2=5.76$ \\
1 & $2(0.6-3)=-4.8$ & $0.6+0.5(0.6-0)-0.1(-4.8)=0.6+0.3+0.48=1.38$ & $(1.38-3)^2=2.6244$ \\
2 & $2(1.38-3)=-3.24$ & $1.38+0.5(1.38-0.6)-0.1(-3.24)=1.38+0.39+0.324=2.094$ & $(2.094-3)^2=0.8196$
\end{tabular}
\end{center}
The iterates rapidly approach the minimiser $w^\star=3$.

%%%%%%%%%%%%%%%%%%%%%%%%%%%%%%%%%%%%%%%%%%%%%%%%%%%%%%%%%%%%
\section*{Assumptions}
%%%%%%%%%%%%%%%%%%%%%%%%%%%%%%%%%%%%%%%%%%%%%%%%%%%%%%%%%%%%
\begin{assumption}[Smoothness]\label{as:smooth}
$f$ is $L$–smooth, i.e., for all $x,y\in\mathbb{R}^d$,
\[
\|\grad f(x)-\grad f(y)\|\le L\|x-y\|,
\]
or equivalently,
\[
f(y)\le f(x)+\langle\grad f(x),y-x\rangle+\frac{L}{2}\|y-x\|^2 .
\tag{A1}
\]
\end{assumption}

\begin{assumption}[Bounded Below]\label{as:bound}
$f$ is bounded from below: $f^\star\triangleq\inf_{x}f(x)>-\infty$.
\tag{A2}
\end{assumption}

\begin{assumption}[Parameter Range]\label{as:param}
The stepsize $\alpha$ and momentum $\beta$ satisfy
\[
0\le\beta<1,\qquad
0<\alpha\le\frac{(1-\beta)^2}{L}.
\tag{A3}
\]
\end{assumption}

No convexity is required; the results hold for any non‑convex $L$‑smooth $f$.

%%%%%%%%%%%%%%%%%%%%%%%%%%%%%%%%%%%%%%%%%%%%%%%%%%%%%%%%%%%%
\section*{Main Convergence Theorem}
%%%%%%%%%%%%%%%%%%%%%%%%%%%%%%%%%%%%%%%%%%%%%%%%%%%%%%%%%%%%
\begin{theorem}[Convergence to a Stationary Point]\label{thm:main}
Assume \ref{as:smooth}–\ref{as:param}.  
Define the Lyapunov sequence
\[
\Phi_k\triangleq f(x_k)-f^\star+\frac{c}{2}\|x_k-x_{k-1}\|^2,
\qquad 
c\triangleq\frac{1-\beta}{\alpha}-\frac{L}{2}.
\]
Then $c>0$ under \ref{as:param}, $\{\Phi_k\}$ is non‑increasing, and for any $K\ge 1$
\[
\frac{1}{K}\sum_{k=0}^{K-1}\|\grad f(x_k)\|^2
\;\le\;
\frac{2\bigl(\Phi_0-\Phi_K\bigr)}{\alpha(1-\beta)} 
\;\le\;
\frac{2\bigl(f(x_0)-f^\star\bigr)}{\alpha(1-\beta)}\frac{1}{K}.
\tag{1}
\]
Consequently,
\[
\lim_{k\to\infty}\|\grad f(x_k)\|=0,
\qquad 
\min_{0\le k\le K-1}\|\grad f(x_k)\|^2 = O\!\left(\frac{1}{K}\right).
\]
\end{theorem}

\section*{Theoretical Properties}
\begin{itemize}
\item \textbf{Stability.} The condition $\alpha\le(1-\beta)^2/L$ guarantees that the Lyapunov function $\Phi_k$ never increases; thus the iterates remain in a bounded sublevel set of $f$.
\item \textbf{Hyperparameter Sensitivity.} Larger $\beta$ accelerates progress when the landscape is locally quadratic, but requires a smaller $\alpha$ (quadratically in $1-\beta$) to preserve stability.
\item \textbf{Comparison.} For $\beta=0$ the method reduces to plain gradient descent and the bound (1) becomes the classical $O(1/K)$ rate for smooth non‑convex functions. Positive $\beta$ does not improve the worst‑case $O(1/K)$ order, but empirically yields faster descent on ill‑conditioned problems.
\end{itemize}

\section*{Discussion}
The theorem shows that, despite the lack of convexity, the Heavy–Ball method converges to a stationary point at the same $O(1/K)$ rate as vanilla gradient descent, provided the stepsize is chosen according to \ref{as:param}. The proof hinges on a carefully crafted Lyapunov function that couples function values and the kinetic energy $\|x_k-x_{k-1}\|^2$. The constant $c$ is positive exactly when the parameter restriction holds, which explains the familiar ``critical stepsize'' $\alpha_{\max}=(1-\beta)^2/L$.

%%%%%%%%%%%%%%%%%%%%%%%%%%%%%%%%%%%%%%%%%%%%%%%%%%%%%%%%%%%%
\section*{Preliminaries (Definitions and Lemmas)}
%%%%%%%%%%%%%%%%%%%%%%%%%%%%%%%%%%%%%%%%%%%%%%%%%%%%%%%%%%%%
\begin{lemma}[Smoothness Inequality]\label{lem:smooth}
For an $L$‑smooth $f$ and any $x,y$,
\[
f(y)\le f(x)+\langle\grad f(x),y-x\rangle+\frac{L}{2}\|y-x\|^2 .
\tag{L1}
\]
\end{lemma}
\begin{proof}
Directly from Assumption~\ref{as:smooth}. \hfill\qedhere
\end{proof}

\begin{lemma}[Descent Lemma for Heavy–Ball]\label{lem:descent}
Let $x_{k+1}$ be generated by \eqref{eq:HB}. For any $c>0$,
\[
\Phi_{k+1}-\Phi_k
\le -\frac{\alpha(1-\beta)}{2}\|\grad f(x_k)\|^2
-\left(c-\frac{L}{2}\right)\|x_{k+1}-x_k\|^2
+\frac{c\beta}{2}\|x_k-x_{k-1}\|^2 .
\tag{L2}
\]
\end{lemma}
\begin{proof}
We expand $\Phi_{k+1}$:
\[
\Phi_{k+1}=f(x_{k+1})-f^\star+\frac{c}{2}\|x_{k+1}-x_k\|^2 .
\]
Apply Lemma~\ref{lem:smooth} with $x=x_k$, $y=x_{k+1}$:
\[
f(x_{k+1})\le f(x_k)+\langle\grad f(x_k),x_{k+1}-x_k\rangle+\frac{L}{2}\|x_{k+1}-x_k\|^2 .
\tag{L2.1}
\]
Insert the update \eqref{eq:HB}:
\[
x_{k+1}-x_k = \beta (x_k-x_{k-1})-\alpha\grad f(x_k).
\tag{L2.2}
\]
Hence
\[
\langle\grad f(x_k),x_{k+1}-x_k\rangle
= \beta\langle\grad f(x_k),x_k-x_{k-1}\rangle-\alpha\|\grad f(x_k)\|^2 .
\tag{L2.3}
\]
Now compute $\Phi_{k+1}-\Phi_k$ using \eqref{L2.1}–\eqref{L2.3} and the definition of $\Phi_k$:
\[
\begin{aligned}
\Phi_{k+1}-\Phi_k
&\le \underbrace{\langle\grad f(x_k),x_{k+1}-x_k\rangle}_{\text{(L2.3)}}
      +\frac{L}{2}\|x_{k+1}-x_k\|^2
      +\frac{c}{2}\|x_{k+1}-x_k\|^2
      -\frac{c}{2}\|x_k-x_{k-1}\|^2 .
\end{aligned}
\]
Group the quadratic terms:
\[
\bigl(L/2+c\bigr)\|x_{k+1}-x_k\|^2
= \left(c-\frac{L}{2}\right)\|x_{k+1}-x_k\|^2
+ L\|x_{k+1}-x_k\|^2 .
\]
Since $L\|x_{k+1}-x_k\|^2\ge0$, discarding it yields an inequality (tightening the bound). Substituting \eqref{L2.2} for $x_{k+1}-x_k$ gives
\[
\|x_{k+1}-x_k\|^2
= \beta^2\|x_k-x_{k-1}\|^2
   -2\alpha\beta\langle\grad f(x_k),x_k-x_{k-1}\rangle
   +\alpha^2\|\grad f(x_k)\|^2 .
\tag{L2.4}
\]
Plug \eqref{L2.4} into the previous inequality and collect terms in $\|\grad f(x_k)\|^2$, $\langle\grad f(x_k),x_k-x_{k-1}\rangle$ and $\|x_k-x_{k-1}\|^2$. After straightforward algebra (see Step~[Step 5] below) we obtain
\[
\Phi_{k+1}-\Phi_k
\le -\frac{\alpha(1-\beta)}{2}\|\grad f(x_k)\|^2
-\left(c-\frac{L}{2}\right)\|x_{k+1}-x_k\|^2
+\frac{c\beta}{2}\|x_k-x_{k-1}\|^2 .
\]
\hfill\qedhere
\end{proof}

%%%%%%%%%%%%%%%%%%%%%%%%%%%%%%%%%%%%%%%%%%%%%%%%%%%%%%%%%%%%
\section*{Main Proof (Step‑by‑Step Derivation)}
%%%%%%%%%%%%%%%%%%%%%%%%%%%%%%%%%%%%%%%%%%%%%%%%%%%%%%%%%%%%
We now prove Theorem~\ref{thm:main}. Every non‑trivial manipulation is labelled.

\subsection*{Step 1 – Choice of the Lyapunov Coefficient}
Define
\[
c\triangleq\frac{1-\beta}{\alpha}-\frac{L}{2}.
\tag{S1}
\]
Assumption~\ref{as:param} implies $\alpha\le (1-\beta)^2/L$, i.e.
\[
\frac{1-\beta}{\alpha}\ge\frac{L}{1-\beta}\ge\frac{L}{2},
\]
hence $c>0$.

\subsection*{Step 2 – Apply Lemma~\ref{lem:descent}}
Using Lemma~\ref{lem:descent} with the $c$ from (S1) yields
\[
\Phi_{k+1}-\Phi_k
\le -\frac{\alpha(1-\beta)}{2}\|\grad f(x_k)\|^2
-\underbrace{\Bigl(c-\frac{L}{2}\Bigr)}_{=\frac{1-\beta}{\alpha}-L}\|x_{k+1}-x_k\|^2
+\frac{c\beta}{2}\|x_k-x_{k-1}\|^2 .
\tag{S2}
\]

\subsection*{Step 3 – Eliminate the Kinetic Term}
Observe that $c-\frac{L}{2} = \frac{1-\beta}{\alpha}-L\ge 0$ by the same bound used in Step~1. Hence we can drop the negative term $-\bigl(c-\frac{L}{2}\bigr)\|x_{k+1}-x_k\|^2$ to obtain a simpler inequality:
\[
\Phi_{k+1}-\Phi_k
\le -\frac{\alpha(1-\beta)}{2}\|\grad f(x_k)\|^2
+\frac{c\beta}{2}\|x_k-x_{k-1}\|^2 .
\tag{S3}
\]

\subsection*{Step 4 – Recursive Summation}
Multiply (S3) by $2/(\alpha(1-\beta))$:
\[
\frac{2}{\alpha(1-\beta)}\bigl(\Phi_{k+1}-\Phi_k\bigr)
\le -\|\grad f(x_k)\|^2
+\frac{c\beta}{\alpha(1-\beta)}\|x_k-x_{k-1}\|^2 .
\tag{S4}
\]
Sum (S4) from $k=0$ to $K-1$:
\[
\frac{2}{\alpha(1-\beta)}\bigl(\Phi_{K}-\Phi_{0}\bigr)
\le -\sum_{k=0}^{K-1}\|\grad f(x_k)\|^2
+\frac{c\beta}{\alpha(1-\beta)}\sum_{k=0}^{K-1}\|x_k-x_{k-1}\|^2 .
\tag{S5}
\]

\subsection*{Step 5 – Bounding the Velocity Sum}
From the definition of $\Phi_k$,
\[
\Phi_k = f(x_k)-f^\star+\frac{c}{2}\|x_k-x_{k-1}\|^2 \ge \frac{c}{2}\|x_k-x_{k-1}\|^2 .
\]
Thus $\|x_k-x_{k-1}\|^2 \le \frac{2}{c}\Phi_k$. Plugging this bound into the rightmost term of (S5) gives
\[
\frac{c\beta}{\alpha(1-\beta)}\sum_{k=0}^{K-1}\|x_k-x_{k-1}\|^2
\le
\frac{c\beta}{\alpha(1-\beta)}\sum_{k=0}^{K-1}\frac{2}{c}\Phi_k
=
\frac{2\beta}{\alpha(1-\beta)}\sum_{k=0}^{K-1}\Phi_k .
\tag{S6}
\]
Because $\Phi_k$ is non‑increasing (see (S3) with the discarded negative term), we have $\Phi_k\le\Phi_0$ for all $k$. Hence
\[
\sum_{k=0}^{K-1}\Phi_k \le K\Phi_0 .
\tag{S7}
\]
Insert (S7) into (S6) and then into (S5):
\[
\frac{2}{\alpha(1-\beta)}\bigl(\Phi_{K}-\Phi_{0}\bigr)
\le -\sum_{k=0}^{K-1}\|\grad f(x_k)\|^2
+\frac{2\beta K}{\alpha(1-\beta)}\Phi_0 .
\tag{S8}
\]

\subsection*{Step 6 – Rearrangement}
Move the gradient sum to the left and divide by $K$:
\[
\frac{1}{K}\sum_{k=0}^{K-1}\|\grad f(x_k)\|^2
\le
\frac{2}{\alpha(1-\beta)}\frac{\Phi_{0}-\Phi_{K}}{K}
+\frac{2\beta}{\alpha(1-\beta)}\Phi_0 .
\tag{S9}
\]
Since $\Phi_K\ge 0$, the first term is bounded by $\frac{2\Phi_0}{\alpha(1-\beta)K}$. Dropping the non‑negative second term (which only enlarges the right‑hand side) yields the clean bound
\[
\frac{1}{K}\sum_{k=0}^{K-1}\|\grad f(x_k)\|^2
\le
\frac{2\bigl(\Phi_0-\Phi_K\bigr)}{\alpha(1-\beta)}\frac{1}{K}
\le
\frac{2\bigl(f(x_0)-f^\star\bigr)}{\alpha(1-\beta)}\frac{1}{K}.
\tag{S10}
\]
Equation~(S10) is exactly the claim (1) of Theorem~\ref{thm:main}. Moreover, because the left‑hand side is a non‑negative average, there must exist at least one index $k\le K-1$ with
\[
\|\grad f(x_k)\|^2 \le \frac{2\bigl(f(x_0)-f^\star\bigr)}{\alpha(1-\beta)}\frac{1}{K},
\]
which gives the $O(1/K)$ rate for the minimal gradient norm. The monotone decreasing property of $\Phi_k$ also implies $\Phi_k\to\Phi_\infty$, and consequently $\|\grad f(x_k)\|\to0$.

\hfill\qed

%%%%%%%%%%%%%%%%%%%%%%%%%%%%%%%%%%%%%%%%%%%%%%%%%%%%%%%%%%%%
\section*{Rate Derivation (With Constants)}
%%%%%%%%%%%%%%%%%%%%%%%%%%%%%%%%%%%%%%%%%%%%%%%%%%%%%%%%%%%%
From Theorem~\ref{thm:main} we have the explicit constant
\[
C\;=\;\frac{2\bigl(f(x_0)-f^\star\bigr)}{\alpha(1-\beta)}.
\]
Thus for any $K\ge1$,
\[
\min_{0\le k\le K-1}\|\grad f(x_k)\|^2 \le \frac{C}{K}.
\]
If one chooses the maximal admissible stepsize $\alpha=(1-\beta)^2/L$, the constant simplifies to
\[
C = \frac{2L\bigl(f(x_0)-f^\star\bigr)}{(1-\beta)^3}.
\]
Hence stronger momentum (larger $\beta$) inflates $C$ polynomially, reflecting the trade‑off between acceleration and stability.

%%%%%%%%%%%%%%%%%%%%%%%%%%%%%%%%%%%%%%%%%%%%%%%%%%%%%%%%%%%%
\section*{Special Cases}
%%%%%%%%%%%%%%%%%%%%%%%%%%%%%%%%%%%%%%%%%%%%%%%%%%%%%%%%%%%%
\begin{itemize}
\item \textbf{No momentum ($\beta=0$).}  The update reduces to $x_{k+1}=x_k-\alpha\grad f(x_k)$, $c=1/\alpha-L/2$, and the bound becomes
\[
\frac{1}{K}\sum_{k=0}^{K-1}\|\grad f(x_k)\|^2\le\frac{2\bigl(f(x_0)-f^\star\bigr)}{\alpha K},
\]
the classic $O(1/K)$ result for gradient descent on smooth non‑convex functions.

\item \textbf{Critical stepsize.}  Setting $\alpha=(1-\beta)^2/L$ yields $c= (1-\beta)/\alpha -L/2 = (1-\beta)L/(1-\beta)^2 - L/2 = L/(1-\beta)-L/2>0$, so the Lyapunov function is well‑defined and the proof remains valid.

\item \textbf{Quadratic case.}  For $f(x)=\frac{1}{2}x^\top Qx$ with $Q\succeq0$, the Heavy‑Ball method coincides with the classical second‑order linear iteration; the same parameter restriction guarantees linear convergence to the unique minimiser.
\end{itemize}

%%%%%%%%%%%%%%%%%%%%%%%%%%%%%%%%%%%%%%%%%%%%%%%%%%%%%%%%%%%%
\section*{Conclusion}
We have presented a complete, step‑by‑step convergence analysis of the Heavy–Ball momentum method for non‑convex $L$‑smooth objectives. Under the simple stepsize condition $\alpha\le(1-\beta)^2/L$, the method enjoys an $O(1/K)$ guarantee on the squared gradient norm, matching the best known rate for first‑order methods without convexity assumptions. The proof relies only on smoothness and a Lyapunov argument, and every algebraic manipulation has been displayed explicitly with line numbers for easy verification.

\end{document}